
\section{Related Work}
\label{sec:related_work}

\paragraph{Classical parallel-queue routing}
JSQ is optimal under identical servers~\cite{weber1978optimal,winston1977optimality}.
JSSQ extends JSQ by directing traffic to the server with the lowest expected sojourn time $(\Qlen{i}+1)/\srvrate_i$ under FCFS~\cite{ephremides1980simple}.
The power-of-$d$-choices method~\cite{mitzenmacher2001power} significantly decreases the tail of the queue-length distribution by sampling $d \geq 2$ servers, although it does not utilize complete queue-state information.
Our focus is on settings where full queue-state information is available and routing decisions can use the entire state vector. Accordingly, the empirical comparisons emphasize full-state baselines such as JSQ and JSSQ rather than partial-information methods.

\paragraph{MaxWeight and Lyapunov-based scheduling}
MaxWeight policies attain optimal throughput in multi-hop constrained queueing networks, where they possess known deterministic stability properties~\cite{tassiulas1992stability}. However, translating a simple MaxWeight assignment rule ($\arg\max_i \srvrate_i \Qlen{i}$) to independent parallel-queue routing behaves differently, as it lacks the structural backpressure mechanisms native to network scheduling.
Neely's drift-plus-penalty framework~\cite{neely2010stochastic} extends Lyapunov methods to stochastic network optimization with utility maximization.
Our proposed framework centers on producing continuous-valued stochastic routing probabilities rather than deterministic threshold assignments like MaxWeight, and uses entropy regularization to actively modulate load balancing. Exploring the exact empirical crossover between approximate softmax MaxWeight routing and JSSQ is deferred.

\paragraph{Restless bandits and Whittle index}
Queue routing can also be framed through the restless-bandit lens. In that setting, Whittle-index policies provide a classical analytical approach~\cite{weber1990index,ninomora2002restless}. GibbsQ differs in both mechanism and objective: it uses a stochastic Gibbs-measure routing law with entropy regularization and develops direct Foster--Lyapunov stability proofs for the resulting CTMCs, rather than deriving an index policy.

\paragraph{Softmax and entropy regularization}
In reinforcement learning, Boltzmann policies serve as standard methods for exploration~\cite{sutton2018reinforcement}. The softmax routing law is equivalent to Boltzmann exploration in multi-armed bandits, where the ``reward'' for server $i$ is the negative queue-length score $-s_i(\Qstate)$. The concentration parameter $\temp$ plays the same role as the inverse temperature in bandit exploration.
The log-sum-exp function acts as the convex conjugate of the negative entropy~\cite{boyd2004convex}.
Such duality connects entropy-regularized routing to the Gibbs-measure perspective adopted here.
Entropy-regularized optimization underpins mirror descent~\cite{nemirovskij1983problem,beck2003mirror}, soft actor-critic~\cite{haarnoja2018soft}, and maximum-entropy RL~\cite{ziebart2010modeling}.
Despite these relationships to the RL and optimization literatures, to our knowledge, explicit Foster--Lyapunov drift analyses for softmax routing in heterogeneous parallel queues have not appeared in the queueing literature.

\paragraph{Foster--Lyapunov and fluid-limit methods for queueing stability}
An alternative stability paradigm uses fluid-limit arguments~\cite{dai1995fluid} commonly adopted in broader queueing networks. GibbsQ adopts the direct Foster--Lyapunov route because it yields explicit, computable drift constants characterizing bounds directly, whereas fluid-limit arguments typically establish stability without constructive bounds. The continuous-time Foster--Lyapunov framework for positive Harris recurrence is the standard direct stability tool for CTMCs~\cite{meyn2009markov}.
Quadratic Lyapunov functions are commonly used in homogeneous settings~\cite{neely2010stochastic}. In heterogeneous settings, weighted analogues are a natural extension because they can encode service asymmetry directly in the drift calculation.
Our UAS proof uses a capacity-weighted Lyapunov function $\lyap_{\mathrm{UAS}} = \frac{1}{2}\sum_i \Qlen{i}^2/\srvrate_i$ together with a prior-weighted KL variational identity, while the calibrated-family certification route studied in Section~\ref{sec:reflected_uas} uses a matched weighted Lyapunov function $\frac{1}{2}\sum_i \Qlen{i}^2/\srvrate_i^{\calbeta}$ and a calibrated prior decomposition. To our knowledge, this combination of proof templates has not appeared in the queueing literature.

\paragraph{Imitation learning and behavior cloning}
Behavior cloning trains a policy through supervised imitation of an expert~\cite{osa2018algorithmic}.
DAgger~\cite{ross2011reduction} corrects the distribution-shift problem of offline BC by iteratively querying the expert under the learner's induced distribution.
Our pipeline uses offline BC for initialization, followed by \reinforce{} fine-tuning to improve upon the cloned policy.
The ablation study uses this setup to isolate the effect of teacher quality.

\paragraph{RL for queueing and resource allocation}
RL has been applied, e.g., to queueing control, scheduling, and routing~\cite{dai2022queueing}.
These approaches generally learn policies end-to-end without analytical stability guarantees.
GibbsQ addresses this gap by offering proven stable analytical policies together with a parameter-conditional certified subset of the reflected family in theory, while the reflected baselines and BC teachers used in the experiments remain empirical unless separately certified.

\paragraph{Mean-field approximations}
For large server pools ($\Nservers \to \infty$), mean-field models provide tractable approximations to queueing dynamics; the supermarket-model literature is a standard example~\cite{vvedenskaya1996queueing,mitzenmacher2001power}.
Our theoretical results are finite-$\Nservers$; stress tests validate UAS behavior up to $\Nservers = 1{,}024$ (Appendix~\ref{sec:appendix_stress}).

\paragraph{Positioning}
GibbsQ occupies a distinct niche at the intersection of these literatures: it combines Foster--Lyapunov stability proofs with entropy-regularized routing and a teacher-guided neural training pipeline, bridging analytical queueing theory and modern learned policies.
